% 一、选题依据

\section{课题来源}

本课题来源于计算机图形学与计算物理仿真领域的前沿研究需求，受到国家自然科学基金和浙江大学研究生科研创新项目的支持。物理模拟的系统性探索可追溯至20世纪中叶，伴随着有限元方法在航空航天与结构工程中的萌芽，人类开始尝试利用严谨的数学工具描述宏观世界的动态演变。进入20世纪80年代末至90年代，物理模拟迎来了一个飞速发展的黄金阶段，计算机辅助工程（CAE）与计算机图形学（CG）两条技术路径并行发展，前者致力于高精度的数值方法还原物理规律，后者则追求模拟的真实感必须植根于底层的物理机制。

在现代计算图形学与计算物理仿真中，弹性体及复杂结构的动力学计算面临着一个普遍且艰巨的挑战："软硬耦合"问题。物理系统在不同尺度或不同组件之间往往展现出截然不同的刚度响应：一维细长杆件中，轴向拉伸刚度（无穷大硬约束）与横向弯曲扭转刚度（极软）交织；三维连续体有限元分析中，局部几何畸变网格会产生巨大的人工数值刚化（Locking）；宏观尺度的复杂机械装备中，子结构交界面处的高频强耦合与内部的低频柔性形变构成典型的软硬耦合现象。在传统的全自由度笛卡尔坐标系或标准节点有限元基底下直接求解这些系统，会导致系统刚度矩阵的条件数急剧爆炸，线性求解器陷入停滞，时间积分步长被极大限制。

本课题旨在通过定制化的基函数与坐标变换方法，在不牺牲物理精确度的前提下突破这一计算瓶颈。核心思想是通过基底的变换重构，将结构性病态转化为易于处理的尺度问题，或直接在降维子空间中规避刚性维度的计算。

\section{课题的研究意义}

物理仿真是计算机图形学、计算力学、机器人学等领域的核心技术，研究成果可应用于工业产品设计、虚拟现实、数字孪生、影视特效等领域。具体而言，我国在航空航天、汽车工业、智能制造等领域对高精度物理仿真有迫切需求，但传统仿真方法在处理复杂耦合系统时存在数值稳定性不足、计算效率低下等问题，主流CAE软件（如ANSYS、Abaqus、Nastran等）主要被欧美企业垄断。本研究提出的基函数重构方法可有效提升仿真系统的鲁棒性和计算效率：通过基函数重构消除系统的结构性病态，显著改善求解器的收敛特性；近线性复杂度的求解算法可大幅缩短仿真计算时间；方法天然支持分布式并行计算，可充分利用国产高性能计算资源。这些成果有望在顶级国际学术会议和期刊（如SIGGRAPH、TOG、TVCG等）发表，为国产CAE软件的数值核心模块提供技术支撑。

\section{国内外研究现状分析}

\subsection{弹性体仿真方法研究现状}

弹性体仿真是计算机图形学和计算力学的核心问题，主要方法包括：质点-弹簧模型计算简单但精度有限；有限元方法具有坚实数学基础但在大变形、网格畸变时面临数值锁死；边界元方法只需边界离散但对非均匀材料处理困难；无网格方法如SPH、MPM避免了网格划分困扰但边界条件处理仍有挑战。近年来研究者们提出了模型降维、自适应网格细化、预处理器优化等多种加速和稳定化技术。

\subsection{细长杆仿真方法研究现状}

细长弹性结构（如毛发、绳索、缝合线等）的仿真是一个具有挑战性的问题。传统节点表示方法需要处理长度约束、四元数单位化约束和轴对齐约束等大量约束条件；Bishop标架方法消除了四元数和对齐约束但单元上的伸长仍作为自由度保留；Super-Helices方法恰好匹配不可伸长Cosserat杆内蕴自由度但积分过程产生稠密质量矩阵；RedMax方法使用块对角预处理器实现高效求解器但对碰撞处理仍有局限。

\subsection{病态问题与条件数优化研究现状}

系统病态性是数值仿真的核心难题，源于矩阵条件数的极端值。现有解决方案包括：预处理器设计（不完全LU分解、多重网格方法、区域分解预处理器）但多为"黑盒"式代数修补；正则化方法（Tikhonov正则化、截断奇异值分解）但可能损失物理精度；基函数优化（谱方法、小波方法、自适应基函数）提供了新思路，但针对软硬耦合问题的专门研究仍较为有限。

\subsection{模态综合法（CMS）研究现状}

模态综合法是解决大规模特征值问题的有效方法，起源于航空航天领域。Craig-Bampton方法将子结构特征模态和界面模态组合构成缩减基，被广泛应用于商业有限元软件；界面模态缩减法通过在Schur补步骤后求解局部特征问题降低计算成本但仍需执行昂贵的Schur补计算；自动多层次子结构法实现了高效的层次化求解。现有CMS方法的主要瓶颈包括：子结构基底无法有效张成目标解空间，Schur补计算代价高昂且内存密集，误差下降缓慢。

\subsection{本研究的切入点与创新点}

针对上述问题，本研究提出"基函数重构"的新范式。理论创新方面，从离散表达层面解决软硬耦合问题，区别于传统的代数预处理器方法。方法创新方面：紧凑坐标表示通过坐标变换消除内在约束，将约束优化转化为无约束优化；相位互补模态基利用多重划分实现基函数相位完备，提升收敛速度；算子自适应基函数识别并过滤高频伪刚性模式，改善系统病态性。应用创新方面，方法可广泛应用于毛发动画、模态分析、有限元仿真等领域。